\documentclass{beamer}
\usepackage[utf8]{inputenc}
\usepackage{color}
\usetheme{default}
\title{DNS Rebinding Attack}
\author{Arseny Kurnikov\\
Hasan Islam\\
Harri Hämäläinen}
\begin{document}


\begin{frame}
\titlepage
\end{frame}

\begin{frame}{Introduction}
    \begin{itemize}
        \item Same origin policy
        \item DNS rebinding attack
        \item DNS rebinding vulnerabilities
        \item Classes of attacks
    \end{itemize}
\end{frame}

\begin{frame}{Same origin policy}
    \begin{itemize}
        \item Origin of a resource can be expressed in a tuple (schema, host, port)
            \begin{itemize}
                \item http://foo.host.example/dir/index.html
                \item http://foo.host.example/anotherDir/index2.html
                \item http://bar.host.example/dir/index.html
                \item https://foo.host.example/dir/index.html
            \end{itemize}
        \item The same origin policy is used to control how resources are loaded
        from one origin and how these resources can interact with each other.
        \item E.g. a client-side script retrieved from an attacker's host
        cannot make requests to another host with different origin
        \begin{itemize}
            \item ... except when it can (CORS, policy files, form POSTs, CSS /
            JavaScript files, ...)
        \end{itemize}
    \end{itemize}
\end{frame}

\begin{frame}{DNS Rebinding Attack}
    DNS rebinding attack exploits the DNS service to circumvent the same origin policies. 
    \begin{enumerate}
        \item The attacker has a web server serving resources from origin X with
        corresponding authoritative name server
        \item A victim retrieves a web page from origin X and
        resolves the hostname to attacker's web server with a DNS query to
        attacker's name server. The DNS response is returned with a very small
        TTL value.
        \item The loaded page contains more requests to be made to the same origin
        (origin X); but new DNS query is required due to small TTL value. This
        time the DNS response is bound to the address of the target web server of
        the attack.
    \end{enumerate}
\end{frame}

\begin{frame}{DNS Rebinding Attack 2}
    \begin{itemize}
        \item DNS rebinding style attacks are known for almost twenty years
        \begin{itemize}
            \item In 1996 DNS rebinding exploit was discovered in JVM
            \footnote{D. Dean, E. W. Felten, and D. S. Wallach. Java
            security: from HotJava to Netscape and beyond. In
            IEEE Symposium on Security and Privacy: Oakland,
            California, May 1996}
        \end{itemize}
        \item Allows two major attack vectors
        \begin{itemize}
            \item \textbf{Firewall circumvention:} 
             to access machines behind firewalls. With direct socket access, the
             attacker can interact with a number of internal services besides
             HTTP.
            \item \textbf{IP hijacking:}
            to access publicly available servers from the client’s IP address.
            This allows the attacker to take advantage of the target’s implicit
            or explicit trust in the client’s IP address.
        \end{itemize}
    \end{itemize}
\end{frame}

\begin{frame}{DNS Rebinding Attack 3}
    \begin{itemize}
        \item Usually the accepted solution is to use \textit{DNS pinning} or
        configure the local DNS subsystem to restrict resolutions to a local
        network block
        \begin{itemize}
            \item \textbf{DNS pinning:} DNS name resolution is cached for a fixed duration, regardless
            of TTL (e.g. in application level caches)
        \end{itemize}
        \item Discussion about DNS rebinding attacks is relevent today as the
        browsers became application platforms supporting all kinds of
        third-party addons introducing new ways to exploit DNS rebinding
        \begin{itemize}
            \item multi-pin vulnerabilities
        \end{itemize}
    \end{itemize}
\end{frame}

\begin{frame}{DNS vulnerabilities}
    \begin{itemize}
        \item \textbf{Time-Varying DNS.}  In 2001, the original attack on Java was extended to use use time-varying DNS:
         \begin{enumerate}
            \item A client visits a malicious web site, attacker.com,containing JavaScript. The attacker’s DNS server is configured to bind attacker.com to the attacker’s IP address with a very short TTL.
 			\item The attacker rebinds attacker.com to the target’s IP address.
			\item The malicious script uses frames or XMLHttpRequest to connect to attacker.com, which now resolves to the IP address of the target’s server.
         \end{enumerate}
         \item \textit{the connection in Step 3 has the same host name as the
         original malicious script, the browser permits the attacker to read the response from the target.
}
    \end{itemize}
\end{frame}


\begin{frame}[fragile]{DNS vulnerabilities}
    \begin{itemize}
        \item \textbf{Flash Player 9.}  Flash Player is vulnerable to the following re-binding attack:
        \begin{itemize}
            \item The client’s web browser visits a malicious web site that embeds a SWF movie.
			\item The SWF movie opens a socket on a port less than 1024 to attacker.com, bound to the attacker’s IP address. Flash Player sends 
            \begin{verbatim}
			<policy-file-request />.
			\end{verbatim}
         \end{itemize}
    \end{itemize}
\end{frame}

\begin{frame}[fragile]{DNS vulnerabilities}
     \begin{itemize}
        \item The attacker responds with the following XML:			        
            \begin{verbatim}
								 <?xml version="1.0"?> 
								 <cross-domain-policy>
								<allow-access-from domain="*" to-ports="*" />
								</cross-domain-policy> 
            \end{verbatim}
        \item The SWF movie opens a socket to an arbitrary port number on attacker.com, which the attacker has rebound to the target’s IP address.
     \end{itemize}
    \textit{The policy XML provided by the attacker in step 3 instructs Flash Player to permit arbitrary socket access to attacker.com.}
\end{frame}

\begin{frame}{Sources}
    \begin{itemize}
        \item Jackson, Collin, et al. "Protecting browsers from DNS rebinding attacks." ACM Transactions on the Web (TWEB) 3.1 (2009): 2.
        \item Ruderman, Jesse, Mozilla developer network, "Same-origin policy",
        \url{https://developer.mozilla.org/en-US/docs/JavaScript/Same_origin_policy_for_JavaScript}
        \item D. Dean, E. W. Felten, and D. S. Wallach. Java security: from
        HotJava to Netscape and beyond. In IEEE Symposium on Security and
        Privacy: Oakland, California, May 1996
    \end{itemize}
\end{frame}

\end{document}
